\documentclass[11pt]{article}
\usepackage{amsmath,amssymb}
\usepackage[margin=1in]{geometry}

\title{The Visible Spectrum Is One Octave}
\author{}
\date{}

\begin{document}
\maketitle

\noindent\textbf{Hook:} The entire range of colors you can see --- from deep red to violet
--- spans almost exactly one octave of electromagnetic frequency.

\medskip

Red light has wavelength $\sim$700\,nm (frequency $\sim$430\,THz). Violet light
has wavelength $\sim$400\,nm (frequency $\sim$750\,THz). The ratio:
$750/430 \approx 1.74$, just shy of $2$.

In music, an octave is a frequency ratio of exactly $2:1$. The visible spectrum
covers $0.81$ octaves --- a major seventh, the interval just below the octave.

This is not a coincidence. The Cayley transform $Q(x) = (1+x)/(1-x)$, which converts
between bounded and unbounded quantities, gives $Q(1/3) = 2$ (the octave) and maps
the reciprocal odd numbers $1/3, 1/5, 1/7, \ldots$ to the superparticular ratios
$2/1, 3/2, 4/3, \ldots$ that define musical consonance in just intonation.

Through the rapidity $\phi = \operatorname{arctanh}(v) = \tfrac{1}{2}\ln Q(v)$,
every ratio becomes a distance:

\begin{center}
\begin{tabular}{lcr}
System & Ratio & Rapidity \\
\hline
Visible spectrum & 700/400\,nm & $0.280$ \\
Musical octave & $2:1$ & $0.347$ \\
Hydrogen Lyman $2\to1$ & $2:1$ & $0.347$ \\
Hydrogen Balmer $3\to2$ & $3:2$ & $0.203$ \\
\end{tabular}
\end{center}

The hydrogen spectral series start on a major chord: the Lyman $2\to 1$
transition is an octave, Balmer $3\to 2$ is a perfect fifth, Paschen $4\to 3$
is a perfect fourth, Brackett $5\to 4$ is a major third. The atomic
energy level ratios $m/n$ are literally the same superparticular ratios
$(n{+}1)/n$ that define musical intervals.

And the entire electromagnetic spectrum --- from radio waves ($\sim10^6$\,Hz)
to gamma rays ($\sim10^{20}$\,Hz) --- spans exactly $46.5$ octaves.
Our eyes see less than one octave of this range. Our ears hear about $10$ octaves
($20$\,Hz to $20$\,kHz). The electromagnetic keyboard has $46$ octaves, and
we see a single key.

\medskip

\noindent\textbf{Rapidity is the universal coordinate.} It measures musical
consonance, relativistic boost, quantum squeezing, information distance,
and tournament coupling strength --- all through the same function
$\operatorname{arctanh}$.

\end{document}
